%%=============================================================================
%% Conclusie
%%=============================================================================

\chapter{Conclusie}%
\label{ch:conclusie}

% TODO: Trek een duidelijke conclusie, in de vorm van een antwoord op de
% onderzoeksvra(a)g(en). Wat was jouw bijdrage aan het onderzoeksdomein en
% hoe biedt dit meerwaarde aan het vakgebied/doelgroep? 
% Reflecteer kritisch over het resultaat. In Engelse teksten wordt deze sectie
% ``Discussion'' genoemd. Had je deze uitkomst verwacht? Zijn er zaken die nog
% niet duidelijk zijn?
% Heeft het onderzoek geleid tot nieuwe vragen die uitnodigen tot verder 
%onderzoek?

\chapter{Conclusie}%
\label{ch:conclusie}

Deze bachelorproef onderzocht hoe een performante provisioner voor virtuele Windows- en Linuxomgevingen kan worden ontworpen en gerealiseerd, die beter inspeelt op de behoeften van het onderwijs dan de huidige Vagrant-gebaseerde aanpak. Het resultaat van dit onderzoek is RustProv, een in Rust geschreven provisioner die virtuele machines in VirtualBox automatisch kan aanmaken, configureren, opstarten en provisionen.

De experimenten tonen aan dat RustProv in de meeste scenario's sneller is dan Vagrant. Bij één Linux-VM is Vagrant in totaal 2,35 seconden sneller, omdat de langere opstarttijd van RustProv wordt gecompenseerd door een aanzienlijk kortere provisioningtijd. Bij drie Linux-VM's is RustProv gemiddeld 74,03 seconden sneller dan Vagrant. Bij één Windows-VM is RustProv in totaal 58,91 seconden sneller, voornamelijk door de lichtere configuratie tijdens het opstarten. Bij drie Windows-VM's bedraagt de tijdswinst gemiddeld 472,10 seconden, omdat RustProv de configuratie van VM's op de achtergrond uitvoert en zo wachttijden vermindert.

Deze resultaten beantwoorden de centrale onderzoeksvraag positief: het is mogelijk om een performante provisioner te ontwerpen die in specifieke scenario's sneller is dan Vagrant, vooral bij het provisionen van meerdere Windows-VM's. RustProv voldoet aan de essentiële functionaliteiten voor een onderwijsprovisioner: het ondersteunt zowel Linux als Windows, biedt configuratie via een leesbaar TOML-bestand en maakt herhaalbare experimenten mogelijk.

De belangrijkste bijdrage van dit werk is het aantonen dat een lichtgewicht, in Rust ontwikkelde provisioner concurrentievoordelen kan bieden ten opzichte van een gevestigde tool zoals Vagrant. RustProv maakt gebruik van de memory safety en performantie van Rust, terwijl het een eenvoudigere architectuur hanteert die gericht is op de kernfunctionaliteit van een provisioner.

Enkele beperkingen moeten worden benadrukt. De metingen werden uitgevoerd op één specifiek testplatform en met een beperkt aantal VM's (één en drie VM's per scenario). De provisioningscripts zijn relatief eenvoudig en vertegenwoordigen niet alle mogelijke onderwijsscenario's. Bovendien is RustProv momenteel beperkt tot VirtualBox als virtualisatieplatform en ondersteunt het enkel NAT- en bridge-netwerkadapters.

Toekomstig onderzoek kan zich richten op het uitbreiden van RustProv met ondersteuning voor additional virtualisatieplatforms zoals Hyper-V of VMware, het toevoegen van geavanceerde netwerkconfiguraties en het verbeteren van de foutafhandeling en logging. Daarnaast kan de impact van RustProv in een reële onderwijscontext worden geëvalueerd door de tool in te zetten in labo-omgevingen en feedback te verzamelen van studenten en docenten.

Samenvattend toont deze bachelorproef aan dat RustProv een levensvatbaar alternatief is voor Vagrant in onderwijsomgevingen, met name wanneer snelheid en herhaalbaarheid belangrijke criteria zijn. De implementatie vormt een basis voor verdere ontwikkeling en validatie in praktische onderwijssettings.